%!TEX root = ../../thesis.tex
% Generated by scripts/evaluate.py -- do not edit; edit the emitter instead.
\begin{figure}[tbp]
  \centering
  \begin{subfigure}[b]{0.48\linewidth}
    \centering
    \begin{tikzpicture}
      \begin{axis}[thesis result plot, width=0.95\linewidth, height=4.0cm,
        xlabel={learner iteration}, ylabel={episode return},
        xmin=0, ymin=0]
      \addplot [draw=none, forget plot, name path=loR]
        table [col sep=comma, x=x, y=lo] {figures/data/input-return-high-and-low.csv};
      \addplot [draw=none, forget plot, name path=hiR]
        table [col sep=comma, x=x, y=hi] {figures/data/input-return-high-and-low.csv};
      \addplot [KITblue15, forget plot] fill between [of=loR and hiR];
      \addplot [KIT line plot A, forget plot]
        table [col sep=comma, x=x, y=mid] {figures/data/input-return-high-and-low.csv};
      \end{axis}
    \end{tikzpicture}
    \caption{high-and-low: return}
  \end{subfigure}
  \hfill
  \begin{subfigure}[b]{0.48\linewidth}
    \centering
    \begin{tikzpicture}
      \begin{axis}[thesis result plot, width=0.95\linewidth, height=4.0cm,
        xlabel={learner iteration}, ylabel={high-coverage byte rate},
        xmin=0, ymin=0, ymax=1]
      \addplot [KIT line plot A, forget plot]
        table [col sep=comma, x=iteration, y=rate] {figures/data/input-rate-high-and-low.csv};
      \addplot [KITblack50, dashed, forget plot] coordinates {(0,0.01553) (99,0.01553)};
      \end{axis}
    \end{tikzpicture}
    \caption{high-and-low: action rate}
  \end{subfigure}

  \vspace{1.4em}

  \begin{subfigure}[b]{0.48\linewidth}
    \centering
    \begin{tikzpicture}
      \begin{axis}[thesis result plot, width=0.95\linewidth, height=4.0cm,
        xlabel={learner iteration}, ylabel={episode return},
        xmin=0, ymin=0]
      \addplot [draw=none, forget plot, name path=loR]
        table [col sep=comma, x=x, y=lo] {figures/data/input-return-sequence.csv};
      \addplot [draw=none, forget plot, name path=hiR]
        table [col sep=comma, x=x, y=hi] {figures/data/input-return-sequence.csv};
      \addplot [KITblue15, forget plot] fill between [of=loR and hiR];
      \addplot [KIT line plot A, forget plot]
        table [col sep=comma, x=x, y=mid] {figures/data/input-return-sequence.csv};
      \end{axis}
    \end{tikzpicture}
    \caption{sequence: return}
  \end{subfigure}
  \hfill
  \begin{subfigure}[b]{0.48\linewidth}
    \centering
    \begin{tikzpicture}
      \begin{axis}[thesis result plot, width=0.95\linewidth, height=4.0cm,
        xlabel={learner iteration}, ylabel={episode length},
        xmin=0, ymin=0]
      \addplot [draw=none, forget plot, name path=loL]
        table [col sep=comma, x=x, y=lo] {figures/data/input-length-sequence.csv};
      \addplot [draw=none, forget plot, name path=hiL]
        table [col sep=comma, x=x, y=hi] {figures/data/input-length-sequence.csv};
      \addplot [KITred15, forget plot] fill between [of=loL and hiL];
      \addplot [KIT line plot B, forget plot]
        table [col sep=comma, x=x, y=mid] {figures/data/input-length-sequence.csv};
      \addplot [KITblack50, dashed, forget plot] coordinates {(0,16) (399,16)};
      \end{axis}
    \end{tikzpicture}
    \caption{sequence: episode length}
  \end{subfigure}

  \vspace{1.4em}

  \begin{subfigure}[b]{0.48\linewidth}
    \centering
    \begin{tikzpicture}
      \begin{axis}[thesis result plot, width=0.95\linewidth, height=4.0cm,
        xlabel={learner iteration}, ylabel={episode return},
        xmin=0, ymin=0]
      \addplot [draw=none, forget plot, name path=loR]
        table [col sep=comma, x=x, y=lo] {figures/data/input-return-open62541.csv};
      \addplot [draw=none, forget plot, name path=hiR]
        table [col sep=comma, x=x, y=hi] {figures/data/input-return-open62541.csv};
      \addplot [KITblue15, forget plot] fill between [of=loR and hiR];
      \addplot [KIT line plot A, forget plot]
        table [col sep=comma, x=x, y=mid] {figures/data/input-return-open62541.csv};
      \end{axis}
    \end{tikzpicture}
    \caption{open62541: return}
  \end{subfigure}
  \hfill
  \begin{subfigure}[b]{0.48\linewidth}
    \centering
    \begin{tikzpicture}
      \begin{axis}[thesis result plot, width=0.95\linewidth, height=4.0cm,
        xlabel={learner iteration}, ylabel={episode length},
        xmin=0, ymin=0]
      \addplot [draw=none, forget plot, name path=loL]
        table [col sep=comma, x=x, y=lo] {figures/data/input-length-open62541.csv};
      \addplot [draw=none, forget plot, name path=hiL]
        table [col sep=comma, x=x, y=hi] {figures/data/input-length-open62541.csv};
      \addplot [KITred15, forget plot] fill between [of=loL and hiL];
      \addplot [KIT line plot B, forget plot]
        table [col sep=comma, x=x, y=mid] {figures/data/input-length-open62541.csv};
      \addplot [KITblack50, dashed, forget plot] coordinates {(0,12) (99,12)};
      \end{axis}
    \end{tikzpicture}
    \caption{open62541: episode length}
  \end{subfigure}
  \caption[Learning on the input-level targets]{
    Learning on the three input-level targets.
    The left columns contains the episode return as a median over the 128 parallel environments
    of one learner iteration as well as the interquartile range.
    The right column presents the return as a program-specific success metric.
    For the high-and-low and sequence targets, the left and right plots are
    identical up to a constant.
    For open62541, there are multiple ways to reach the same return, so the figures
    don't align. }
  \label{fig:input-learning}
\end{figure}
